\section{引言：Vibe Coding 时代的质量困境}

在 AI 辅助编程（Vibe Coding）快速普及的今天，开发者面临一个核心问题：\textbf{如何在使用 Codex、Claude Code 等纯自然语言编程工具时，依然保证代码的准确性与可维护性}。本期围绕开源项目 \textbf{Superpowers}（GitHub Star 7 万+）这一组高质量 Skills，系统性地回答这一问题。

\begin{knowledgebox}{Vibe Coding 的工具光谱}
讲者将当前 AI 编程工具分为三个层级：
\begin{enumerate}
    \item \textbf{VS Code（无 AI 插件）}：纯粹的代码编辑器 / IDE。
    \item \textbf{Cursor}：带有 Tab 模型（大规模 on-policy RL 训练）的轻量级 AI 辅助，支持对话式开发与 diff GUI。
    \item \textbf{Codex / Claude Code}：纯自然语言编程 CLI，无 IDE，只需描述需求并验证输出。
\end{enumerate}
\end{knowledgebox}

传统方法中，代码质量通过 \textbf{TDD}（测试驱动开发）、\textbf{Code Review}、\textbf{Pre-commit Hook}（如 Ruff 静态检查）等流程来保障。在 Vibe Coding 时代，Superpowers 将这些顶级软件工程实践"硬编码"为 AI Coding Agent 的规范操作系统——通过 Skill 指令而非 Hook 脚本，强制 Agent 克服大模型常见的幻觉、懒惰和讨好倾向。

\begin{importantbox}{核心理念："慢其实快"}
不加任何 Skill、不加任何约束，模型可能很快就产出代码，但代码不准、不可维护。有了 Superpowers 的约束后，虽然每一步都更慎重，但代码质量有保障，长期迭代效率更高——这就是"慢其实快"。
\end{importantbox}


